\section*{Bab \ref{ch:g1:the-five-senses} --- Lima Indra}

\begin{solution}{exo:g1:the-five-senses:1}
Penglihatan --- \omterm{def:g1:the-five-senses:sense}{mata}; pendengaran --- \omterm{def:g1:the-five-senses:sense}{telinga}; penciuman --- hidung;
pengecapan --- lidah; perabaan --- \omterm{def:g1:the-five-senses:sense}{kulit}.
\end{solution}

\begin{solution}{exo:g1:the-five-senses:2}
(a) \omterm{def:g1:the-five-senses:sense}{Matanya}. (b) \omterm{def:g1:the-five-senses:sense}{Kulitnya}. (c) Hidungnya. (d) \omterm{def:g1:the-five-senses:sense}{Telinganya}.
\end{solution}

\begin{solution}{exo:g1:the-five-senses:3}
\omterm{def:g1:the-five-senses:sense}{Kulit}, alat perabaan. Ia menyelimuti seluruh tubuh, dari ujung kepala
sampai ujung kaki.
\end{solution}

\begin{solution}{exo:g1:the-five-senses:4}
\omterm{def:g1:the-five-senses:sense}{Kulit} tangan mengirim pesannya (dingin, keras, bergerigi); otak
memahaminya lalu menyebut nama bendanya.
\end{solution}

\begin{solution}{exo:g1:the-five-senses:5}
Pengecapan dan penciuman. Hidung mengerjakan separuh dari cita
rasanya, dan karena itu makanan terasa hambar ketika hidung tersumbat.
\end{solution}

\begin{solution}{exo:g1:the-five-senses:6}
Misalnya: jangan pernah memandang lurus ke Matahari --- melindungi
\omterm{def:g1:the-five-senses:sense}{mata}; jagalah bunyi tetap lembut --- melindungi \omterm{def:g1:the-five-senses:sense}{telinga}; tiuplah
makanan panas --- melindungi lidah; jangan memasukkan apa pun ke
\omterm{def:g1:the-five-senses:sense}{telinga} atau hidung --- melindungi \omterm{def:g1:the-five-senses:sense}{telinga} dan hidung; cahaya yang
cukup untuk membaca --- melindungi \omterm{def:g1:the-five-senses:sense}{mata}.
\end{solution}

\begin{solution}{exo:g1:the-five-senses:7}
Karena bunyi yang terlalu keras melukai \omterm{def:g1:the-five-senses:sense}{telinga}, dan \omterm{def:g1:the-five-senses:sense}{telinga} yang
terluka mungkin tetap terluka. Penutupnya menjaga bunyi mesin tetap
cukup lembut bagi pendengaran seumur hidup.
\end{solution}

\begin{solution}{exo:g1:the-five-senses:8}
Jawaban terbuka --- kebanyakan orang menyebut dua atau tiga dari tiga.
Hidung mengenali makanan yang sudah dikenalnya dengan baik, bahkan
tanpa bantuan \omterm{def:g1:the-five-senses:sense}{mata}.
\end{solution}

\begin{solution}{exo:g1:the-five-senses:9}
\omterm{def:g1:the-five-senses:sense}{Mata} hanya \emph{melihat} --- ia mengumpulkan cahaya lalu mengirim
sebuah pesan. Otaklah yang memahami: ia menerima pesannya lalu berkata
``itu Nenek'', ``itu sekolahku''. Indra mengumpulkan; otak memahami.
\end{solution}

\begin{solution}{exo:g1:the-five-senses:10}
Membaca: dengan ujung jari --- perabaan mengambil alih penglihatan,
merasakan titik-titik timbul di halamannya. Menyeberang jalan: dengan
\omterm{def:g1:the-five-senses:sense}{telinga} --- mendengarkan mobil yang datang dan pergi --- sering
disertai tongkat yang ujungnya membiarkan perabaan memeriksa tanah di
depannya.
\end{solution}
